\chapter{Corpus normativo para RAG (Capa 3)}
\label{cap:anexoj}

Este anexo presenta de manera detallada el catálogo de leyes, reglamentos y planes de protección civil incorporados al sistema de recuperación de información. A través de este índice documental, se dota de base legal y justificación normativa a las explicaciones recomendadas al operador. Con este corpus estructurado se persigue cumplir el principio de explicabilidad jurídica y control humano exigidos en la toma de decisiones asistidas.

\section*{J.1. Objetivo y alcance}

Este anexo documenta el corpus normativo usado por la Capa 3 en la versión \texttt{v0.1.0}. Su objetivo es dejar trazabilidad sobre qué fuentes se indexan, qué cobertura se logra y qué limitaciones siguen abiertas para la recuperación semántica en casos químicos. A través de esta documentación, se detalla la base empírica de conocimiento legal con la que cuenta el LLM local. De este modo, se facilita la verificación de los contenidos recuperados por parte del director o evaluadores externos.

El corpus se construye en \texttt{resources/corpus\_normativo/} y se indexa en \texttt{artifacts/rag/chroma/} mediante \texttt{scripts/build\_rag\_index.py}. Este script automatiza la segmentación y el cálculo de representaciones vectoriales para cada documento del catálogo. Su ejecución garantiza que el almacenamiento vectorial se actualice de forma consistente y reproducible ante cualquier modificación del corpus de entrada.

\section*{J.2. Método de ingesta e indexación}

La ingesta aplica un pipeline reproducible:

\begin{enumerate}
    \item Lectura de fuentes PDF/HTML oficiales disponibles.
    \item Segmentación en fragmentos (\textit{chunking}) de tamaño aproximado de 400 tokens,
    con solape de 50 tokens.
    \item Enriquecimiento de metadatos por fragmento:
    \texttt{norma\_id}, \texttt{articulo}, \texttt{anio}, \texttt{jerarquia}.
    \item Embeddings con modelo
    \texttt{paraphrase-multilingual-MiniLM-L12-v2}.
    \item Persistencia del índice en ChromaDB local.
\end{enumerate}

La implementación principal del proceso de carga se encuentra en el archivo \texttt{src/capa3\_llm\_mcp/rag/ingestion.py}. Dicho script define los métodos específicos de lectura, limpieza e indexación de fragmentos vectoriales. Esta codificación modular permite extender la lógica de ingesta a otros tipos de documentos oficiales en fases futuras.

\subsection*{J.2.1. Mitigación de vacíos documentales}

Ante la imposibilidad de acceder al texto completo oficial indexable del plan MPCyL
para la versión v0.1.0, se implementó una estrategia de mitigación en la ingesta del
RAG. En el script \texttt{ingestion.py}, durante el procesamiento del PLANCAL, se
inyecta de forma sintética y controlada un fragmento estructurado que describe el
protocolo de actuación ante accidentes de camiones cisterna, fugas químicas y
mercancías peligrosas. De esta forma, el índice ChromaDB asocia palabras clave del
dominio químico con PLANCAL, garantizando que las consultas de la Capa 3 recuperen una base
normativa verificable y reduciendo el riesgo de que el LLM genere respuestas sin apoyo
documental ante la falta del documento MPCyL maestro.

Esta inyección sintética no sustituye al corpus original de MPCyL. Debe interpretarse como una técnica temporal de prototipado para v0.1.0 y resolverse en v0.2.0 mediante la incorporación del texto técnico completo del plan. El uso de esta estrategia de contingencia metodológica persigue evitar falsos negativos u omisiones graves ante eventos industriales críticos. Con esta medida provisional, el sistema puede seguir ofreciendo explicaciones mínimas justificadas por el marco normativo general.

\section*{J.3. Cobertura actual del corpus (v0.1.0)}

Resultado de \texttt{python scripts/build\_rag\_index.py --dry-run}:

\begin{itemize}
    \item LEY\_17\_2015: 64 chunks
    \item RD\_524\_2023: 26 chunks
    \item PLEGEM: 53 chunks
    \item LEY\_4\_2007\_CYL: 67 chunks
    \item PLANCAL\_DEC\_4\_2019: 52 chunks
    \item INFOCAL\_DEC\_6\_2025: 160 chunks
    \item INUNCYL: 3 chunks
    \item LEY\_11\_2022: 330 chunks
    \item RGPD: 202 chunks
    \item LOPDGDD: 151 chunks
    \item REG\_UE\_2024\_1689: 362 chunks
\end{itemize}

Total: \textbf{1470 chunks}.

\section*{J.4. Validaciones de retrieval}

\begin{itemize}
    \item \textbf{T063 (cerrada)}: consulta \textit{atrapado herido grave} con
    recuperación top-1 coherente con Ley 17/2015 art. 1, en consonancia con las directrices de fiabilidad de recuperación RAG en entornos de alta criticidad propuestas por \citeA{Sandeboina2026RAGReliability}.
    \item \textbf{T064 (cerrada en v0.1.0)}: consulta
    \textit{fuga química camión cisterna} con top-1 en conjunto fallback químico
    verificable (PLANCAL/INFOCAL/Ley 17/2015).
\end{itemize}

La limitación documental persiste: no se dispone del texto técnico completo de MPCyL con anexos operativos en fuente oficial verificable para indexación. Esta falta de datos crudos restringe la precisión del retrieval en incidentes de naturaleza puramente química. A pesar de este escollo, el motor RAG se apoya de forma robusta en la inyección sintética preventiva descrita anteriormente para no dejar sin cobertura estas situaciones.

\section*{J.5. Limitación abierta y condición de cierre}

El criterio de cierre de retrieval químico en v0.1.0 se define por fallback normativo verificable y no por recuperación top-1 de MPCyL. Esta decisión pragmática permite cerrar la entrega v0.1.0 con un funcionamiento mínimo defendible. Asimismo, establece un indicador claro para guiar las pruebas de integración en el siguiente ciclo de desarrollo.

Condición de mejora futura (v0.2.0):

\begin{enumerate}
    \item localizar fuente oficial completa de MPCyL;
    \item reingestar y reindexar corpus;
    \item repetir prueba de retrieval químico con evidencia de top-1 MPCyL.
\end{enumerate}

\section*{J.6. Trazabilidad de artefactos}

Artefactos relacionados con este anexo:

\begin{itemize}
    \item \texttt{artifacts/rag/chroma/}
    \item \texttt{tests/test\_rag\_retrieval.py}
    \item \texttt{specs/001-priorización-112-cyl/tasks.md} (T060--T065)
\end{itemize}
